%% ─────────────────────────────────────────────────────────────────────────────
%%  El espacio BV([a,b],X) como ENO
%% ─────────────────────────────────────────────────────────────────────────────
\section{El espacio $\BVX$}

\begin{frame}{}
\begin{block}{}
\centering\large
El espacio $BV([a,b],X)$
\end{block}
\end{frame}

%──────────────────────────────────────────────────────────────────────────────

\begin{frame}{Definición y estructura de Banach}
\framesubtitle{$BV([a,b],X)$ hereda la completitud de $X$}

\begin{definition}
Sea $(X,\|\cdot\|,\Xp)$ un ENO. Se define
\[
\BVX = \bigl\{f:[a,b]\to X : \Var_a^b(f,X)<\infty\bigr\}
\]
con norma $\|f\|_{\BVX} = \|f(a)\| + \Var_a^b(f,X)$.
\end{definition}

\medskip

\begin{theorem}
\begin{enumerate}
  \item Para toda $f\in\BVX$,
        $\|f\|_\infty \le \|f\|_{\BVX}$.
  \item $(\BVX,\|\cdot\|_{\BVX})$ es un \textbf{espacio de Banach}.
\end{enumerate}
\end{theorem}

\end{frame}

%──────────────────────────────────────────────────────────────────────────────

\begin{frame}{$\BVX$ como espacio normado ordenado}
\framesubtitle{El orden se hereda del cono de funciones crecientes}

Sea $P_c = \{f\in\BVX : f \text{ creciente},\ f\ge 0\}$.

\medskip

\begin{theorem}[\cite{perez2022jordan}]
\begin{itemize}
  \item[--] $P_c$ es un cono para $(\BVX,\|\cdot\|_{\BVX})$.
  \item[--] Si $\Xp$ es \textbf{sólido}, entonces $P_c$ es cono
            \textbf{generador} de $\BVX$.
  \item[--] Si $P_c$ es cono generador de $\BVX$, entonces $\Xp$
            es generador.
\end{itemize}
\end{theorem}

\end{frame}

%──────────────────────────────────────────────────────────────────────────────

\begin{frame}{Propiedades del cono de $\BVX$}
\framesubtitle{Normalidad y transferencia de estructura}

\begin{block}{Normalidad}
Si $\Xp$ es \textbf{normal}, entonces $P_c$ es cono \textbf{normal}
en $\BVX$.
\end{block}

\medskip

\begin{alertblock}{Lectura}
Las propiedades del cono se transfieren: $\BVX$ es un ENO y
sus propiedades reflejan las del cono de $X$.
\end{alertblock}

\end{frame}

%──────────────────────────────────────────────────────────────────────────────

\begin{frame}{Estructura algebraica de $\BVX$}
\framesubtitle{Cuando $X$ es álgebra}

\begin{theorem}[\cite{perez2022jordan}]
Si $X$ es un álgebra con $\|fg\|\le\|f\|\,\|g\|$,
entonces $(\BVX,3\|\cdot\|_{\BVX})$ también es un \textbf{álgebra}.
\end{theorem}

\medskip

\begin{alertblock}{Importancia}
Esta estructura permitirá aplicar teoremas de punto fijo en $\BVX$
como espacio normado ordenado con cono y álgebra.
\end{alertblock}

% TODO: agregar referencia exacta para el álgebra BV.

\end{frame}
